\documentclass[11pt,a4paper]{article}

% Encodage et langue
\usepackage[T1]{fontenc}
\usepackage[utf8]{inputenc}
\usepackage[french]{babel}
\usepackage{lmodern}

% Mise en page et mathématiques
\usepackage[a4paper,margin=2.3cm]{geometry}
\usepackage{amsmath,amssymb,mathtools}
\usepackage{enumitem}
\usepackage{microtype}
\usepackage{xcolor}
\usepackage{hyperref}

\hypersetup{
    colorlinks=true,
    linkcolor=blue,
    urlcolor=blue,
    pdftitle={Évaluation TP -- Classification et filtrage bayésien},
    pdfauthor={RT4 -- INSAT}
}

\setlength{\parskip}{0.45em}
\setlength{\parindent}{0pt}
\setlist[itemize]{leftmargin=1.5em,itemsep=0.35em}
\setlist[enumerate]{leftmargin=1.8em,itemsep=0.6em}

\newcommand{\R}{\mathbb{R}}
\newcommand{\E}{\mathbb{E}}
\newcommand{\CC}{\mathbf{CC}}
\newcommand{\D}{\mathbf{D}}
\newcommand{\xv}{\mathbf{x}}
\newcommand{\wv}{\mathbf{w}}
\newcommand{\cv}{\mathbf{c}}
\newcommand{\rv}{\mathbf{r}}

\begin{document}

\begin{center}
    {\Large\bfseries Évaluation TP : Classification et filtrage bayésien}\\[0.35em]
    {\large RT4 -- INSAT}\\
    Responsable : R. Amara
\end{center}

\vspace{0.5em}
\hrule
\vspace{0.8em}

Dans le cours de classification bayésienne, on a pu déterminer l'estimateur
MAP (\textit{Maximum A Posteriori}) du signal de référence $d(n)$,
connaissant le vecteur $\xv(n)$ (le régresseur), pour l'application de
déconvolution pour laquelle
\[
    d(n)\in\{-1,+1\}.
\]

L'observation est donnée par
\[
    x(n)=\CC * D_M(n)+v(n), \qquad n\in\{1,\ldots,N\},
\]
où $v(n)$ est un bruit blanc gaussien :
\[
    v(n)\sim\mathcal{N}(0,\sigma_v^2).
\]
En développant la convolution,
\[
    x(n)
    =\sum_{i=0}^{M-1}c(i)d(n-i)+v(n)
    =CC^{T} D_M(n)+v(n).
\]

Ainsi, le vecteur régresseur est
\[
    \xv(n)=
    \begin{bmatrix}
        x(n)\\
        x(n-1)\\
        \vdots\\
        x(n-P+1)
    \end{bmatrix}.
\]

La matrice de convolution $\CC$ est la matrice circulante donnée par
\[
\CC=
\begin{bmatrix}
c(0) & \cdots & c(M-1) & 0 & \cdots & 0\\
0 & c(0) & \cdots & c(M-1) & \cdots & 0\\
\vdots & \ddots & \ddots & & \ddots & \vdots\\
0 & \cdots & 0 & c(0) & \cdots & c(M-1)
\end{bmatrix}.
\]

$\D_M(n)$ est le vecteur régresseur construit à partir des symboles, tel que
\[
    \D_M(n)=
    \begin{bmatrix}
        d(n) & d(n-1) & \cdots & d(n-P-M+2)
    \end{bmatrix}^{T}.
\]

\section{Classification bayésienne}

\begin{enumerate}[label=\arabic*.]

\item \textbf{Affichage des deux classes auxquelles appartiennent les observations.}

Afficher les deux classes auxquelles appartiennent $\xv(n)$ et relatives à
$d(n)=+1$ et à $d(n)=-1$, pour $M=2$ et $P=2$.

\begin{itemize}
    \item Pour $c(0)=1$ et $c(1)=0.5$, déterminer la matrice $\CC$,
    les vecteurs non bruités $\xv(n)$, à savoir
    \[
        \xv_i^{+}=\CC\,\D_i^{+},
        \qquad
        \xv_i^{-}=\CC\,\D_i^{-},
    \]
    avec $\D_i^{+}$ les valeurs prises par $\D_M(n)$ lorsque $d(n)=+1$,
    et $\D_i^{-}$ les valeurs prises par $\D_M(n)$ lorsque $d(n)=-1$.

    L'implémentation à réaliser doit accepter n'importe quelle valeur pour
    les coefficients $\cv$. Le paramètre $M$ est la longueur du vecteur
    $\cv$. La valeur de $P$ sera fixée au début du code, ainsi que le nombre
    d'échantillons à traiter $N$.

    Les motifs de la classe $\mathcal{E}^{+}$ seront représentés par
    le symbole « + » et ceux de $\mathcal{E}^{-}$ par le symbole « o ».

    \item Après avoir affiché les motifs qui représentent les données
    non bruitées de $\mathcal{E}^{+}$ et $\mathcal{E}^{-}$, superposer
    les valeurs prises par $\xv(n)$ en présence de bruit avec le motif
    « * », avec une autre couleur. Vous pouvez tester des valeurs de
    $\sigma_v^2$ pour lesquelles les classes sont linéairement séparables
    ou non.

    \item Implémenter le calcul du vecteur de probabilités
    \[
        \mathbf{P}(n)=
        \begin{bmatrix}
            P_{+}(n)\\
            P_{-}(n)
        \end{bmatrix},
    \]
    avec $P_{+}(n)$ et $P_{-}(n)$ les probabilités que $\xv(n)$
    appartienne respectivement à $\mathcal{E}^{+}$ et $\mathcal{E}^{-}$.
    On rappellera les expressions de $P_{+}(n)$ et $P_{-}(n)$.

    \item Déterminer le taux d'erreur binaire (TEB), qui est le même que
    le taux d'erreur de classification pour
    \[
        \mathrm{TEB}=\frac{\text{nombre d'erreurs}}{N},
    \]
    où
    \[
        \mathrm{erreur}(n)=
        \begin{cases}
            1, & \text{si } \widehat d_{\mathrm{MAP}}(n)\neq d(n),\\
            0, & \text{sinon}.
        \end{cases}
    \]
    L'estimateur MAP est défini par
    \[
        \widehat d_{\mathrm{MAP}}(n)=
        \begin{cases}
            +1, & \text{si }P_{+}(n)>P_{-}(n),\\
            -1, & \text{si }P_{+}(n)<P_{-}(n).
        \end{cases}
    \]

    \item Afficher la frontière de décision (\textit{decision boundary})
    \[
        \partial\mathcal{E}
        =
        \left\{
        \xv\in\text{espace de }\xv(n)
        \;\middle|\;
        P_{+}(\xv)=P_{-}(\xv)
        \right\}.
    \]
    Il faudra discrétiser l'espace de $\xv(n)$ (autour des classes
    $\mathcal{E}^{+}$ et $\mathcal{E}^{-}$) et afficher l'ensemble des
    points $\xv$ pour lesquels
    \[
        \left|P_{+}(\xv)-P_{-}(\xv)\right|<\varepsilon.
    \]
    À cause de la précision numérique, on ne pourra pas avoir
    $P_{+}=P_{-}$ exactement ; il faut donc régler adéquatement
    $\varepsilon$.

    \item Implémenter de même les coefficients du filtre de Wiener
    $\wv_{\mathrm{opt}}$ (tels que cela a été fait en cours/TP), puis
    afficher l'hyperplan de Wiener séparateur des deux classes
    $\mathcal{E}^{+}$ et $\mathcal{E}^{-}$, défini par
    \[
        H_{\mathrm{Wiener}}^{\mathrm{HP}}
        =
        \left\{
        \xv\;\middle|\;
        \wv_{\mathrm{opt}}^{T}\xv=0
        \right\}.
    \]
    Interpréter l'optimalité du filtre de Wiener et celle du classifieur
    bayésien implémenté, en faisant varier $\sigma_v^2$.

    \item Pour une valeur de $\sigma_v^2$ assurant la séparabilité
    linéaire de $\mathcal{E}^{+}$ et $\mathcal{E}^{-}$, proposer une
    méthode afin de trouver un meilleur hyperplan séparateur que celui
    de Wiener. Implémenter votre méthode et comparer les TEB.
    Il est préférable que la méthode soit issue de votre propre réflexion.
    Un essai de quantification du degré de séparabilité linéaire des deux
    classes peut être envisagé (penser à voir le problème géométriquement).

    \item Changer la valeur de $\cv$ et observer de même la disposition
    des classes, le degré de séparabilité linéaire et celui obtenu par
    classification bayésienne. On pourra présenter le TEB en fonction
    du RSB (rapport signal à bruit) défini par rapport à $\xv(n)$ par
    \[
        \mathrm{RSB}
        =
        10\log_{10}
        \left(
        \frac{\cv^{T}\mathbf{R}_{d}\cv}{\sigma_v^2}
        \right).
    \]
    Pour le classifieur bayésien, le filtre de Wiener mis en œuvre pour
    la déconvolution, ainsi que le filtre implémentant le nouvel
    hyperplan séparateur que vous avez proposé, seront également étudiés.

    Pour le choix de $\cv$, montrer en dimension $M=3$ (filtre
    convolutionnel à 3 coefficients : $c(0)$, $c(1)$ et $c(2)$).

    Pour $M>3$, utiliser un outil de visualisation de dimension tel que
    le t-SNE afin d'afficher les éléments des classes.
\end{itemize}

\end{enumerate}

\section{Étude de l'architecture DNN : DiffWave}

\begin{enumerate}[label=\alph*.]

\item \textbf{Identification de l'architecture DiffWave et de ses applications.}

Identifier l'architecture DiffWave et les applications pour lesquelles
elle est dédiée. Pour chaque application, indiquer ce qui constitue
l'entrée et la sortie de l'architecture DiffWave.

Expliquer le principe de l'approche d'apprentissage de DiffWave en
utilisant des illustrations et des équations. Qu'entend-on par les
opérations de diffusion et de conditionnement ?

Quelles sont les différences avec les vocodeurs tels que WaveNet ?

\item \textbf{Présentation de l'architecture.}

Présenter l'architecture avec ses différents modules. Chaque module est
à détailler (structure, entrée/sortie, rôle).

Accorder une attention particulière au bloc résiduel à base de
convolution dilatée.

\item \textbf{Critère de coût et ressources.}

Quel est le critère ou la fonction coût utilisée pour l'apprentissage ?
S'instruire à propos de la latence et du « memory usage » de DiffWave.

\item \textbf{Identification de versions pré-entraînées et évaluation.}

Identifier quelques versions pré-entraînées de DiffWave que vous pouvez
tester (donner les sources et les dépôts).

Prévoir le test d'un exemple de modèle DiffWave pré-entraîné, sur un
dataset (ou une partie d'un dataset) de votre choix. On montrera la
performance du modèle pour l'application choisie (speech enhancement,
denoising, séparation de sources, etc.) en termes d'indicateurs adéquats.

Vous pouvez toujours construire votre propre scénario de test. Par exemple :
\begin{itemize}
    \item Bruiter un dataset de signaux audio (issus d'enregistrements
    que vous avez réalisés), puis appliquer DiffWave pour les débruiter.
    \item Envisager d'utiliser DiffWave pour résoudre un problème de type
    \textit{cocktail party}, afin d'améliorer la qualité d'un signal audio
    noyé dans une ambiance sonore particulière (ambiance d'une zone
    urbaine, gare, etc.). Commencer donc par créer le mélange de signaux
    audio.
\end{itemize}

\item \textbf{Proposition d'un cas d'utilisation (\textit{use case}).}

Proposer un \textit{use case} de DiffWave afin de générer des signaux de
parole qui pourront être prononcés par un malentendant. Vous commencerez
par décrire votre démarche afin d'implémenter un tel module en détaillant
les étapes.

\end{enumerate}

\vspace{1em}
\begin{center}
    \textit{La compréhension du modèle DNN DiffWave et de son fonctionnement
    sera fortement valorisée.}
\end{center}

\end{document}